	\section{Introduction}
	\label{sec:intro}
	
	Matroids are simple combinatorial structures, providing
	a unified abstraction for independence systems such as linear independence in a vector space, or cycle-free subsets of edges in a given graph. A {\em matroid} is a set system $(E, \cI)$, where $E$ is a finite set and $\cI \subseteq 2^E$ are the {\em independent sets (IS)} such that $(i)$ $\emptyset \in \cI$, $(ii)$ for all $A \in \cI$ and $B \subseteq A$ it holds that $B \in \cI$, and $(iii)$ for all $A,B \in \cI$ where $|A| > |B|$, there is $e \in A \setminus B$ such that $B \cup \{e\} \in \cI$.\footnote{Properties $(ii)$ and $(iii)$ are known as {\em hereditary property},
		and  {\em exchange property}.}
	
	While serving as a generic abstraction for numerous applications, matroids possess useful combinatorial properties that allow the development of efficient algorithms. 
	These algorithms include such canonical results as the classic greedy approach for finding a maximum weight independent set of a matroid (see, e.g., \cite{CLRS22}), Edmond's algorithm for matroid partitioning \cite{edmonds1965minimum}, and Lawler's algorithm for matroid intersection \cite{La75}. In all of the above, polynomial running time is enabled due to the structure of the problem $-$ a single objective function with a matroid constraint. However, in many natural applications, there is an added {\em linear} constraint. 
	
	Consider, for example, the problem of finding a maximum independent set in a matroid subject to a budget constraint. Formally, we are given 
	a set of elements $E$, a membership oracle for a collection of independent sets $\II \subseteq 2^{E}$ of a matroid $(E, \II)$, %where $(E, \II)$ is a matroid, 
	a budget $L>0$, a weight function~$w:E\to \mathbb{R}_{\geq 0}$, and a value function $v:E\rightarrow \mathbb{R}_{\geq 0}$. A {\em solution} for the problem is an independent set $S \in \II$ of total weight at most $L$, i.e., $w(S) \leq L$.\footnote{For every set $X$, a function $f:X \rightarrow \mathbb{R}_{\geq 0}$ and $Y \subseteq X$ we define $f(Y) = \sum_{e \in Y} f(e)$.} The {\em value} of a solution $S$ is given by $v(S)$,
	and the objective is to find a solution of maximum value. This problem, known as {\sc budgeted matroid independent set (BM)}, is a generalization of the classic $0/1${\sc-knapsack}, which is NP-hard
	and therefore unlikely to admit an exact polynomial-time algorithm. Thus, obtaining efficient approximations has been a main focus in the study of BM.

	
 For an instance $I$ of an optimization problem $\cG$, let $\OPT_{\cG}(I)$ be the value of an optimal solution for $I$. For some $\rho \geq 1$, a {\em $\rho$-approximate solution} $S$ for $I$ is a solution of value $v \geq \frac{\OPT_{\cG}(I)}{\rho}$ ($v \leq \rho \cdot \OPT_{\cG}(I)$) if $\cG$ is a maximization (minimization) problem. 
 We say that $\cA$ is a {\em randomized $\rho$-approximation algorithm} for $\cG$ if,
given an instance $I$ of $\cG$, $\cA$
returns with probability at least $\frac{1}{2}$ a $\rho$-approximate solution for $I$ $-$ if a solution exists. If no solution exists $-$ $\cA$ returns that $I$ does not have a solution.
	

	
	
	
	Let $|I|$ be the encoding size of an instance~$I$ of a problem $\cG$. 
	A (randomized) {\em polynomial-time approximation scheme} (PTAS)
	for $\cG$ is a family of algorithms $(A_{\eps})_{\eps>0}$ such that, for any $\eps>0$, $A_{\eps}$ is a (randomized) polynomial-time $(1+ \eps)$-approximation algorithm for $\cG$. 
	A (randomized) {\em Efficient PTAS} (EPTAS) is a (randomized) PTAS $(A_{\eps})_{\eps>0}$ with running time of the form $f\left(\frac{1}{\eps}\right) \cdot |I|^{O(1)}$, where $f$ is an arbitrary computable function. %We say that 	
	The strong dependence of run-times on the error parameter, $\eps > 0$, often renders the above schemes highly impractical. This led to the study of the following more desirable class of schemes. A (randomized) approximation scheme 
	$(A_{\eps})_{\eps>0}$ is a (randomized) {\em Fully PTAS} (FPTAS) if the running time of $A_{\eps}$ is of the form $ {\left(\frac{|I|}{\eps}\right)}^{O(1)}$.\footnote{To better distinguish between EPTAS and FPTAS, we use throughout the paper Efficient PTAS and Fully PTAS.}
	
In the past decades, BM was shown to admit a PTAS~\cite{BBGS11,CVZ11,GZ10}, and more recently an Efficient PTAS \cite{DKS23,DKS23b}. As the special case of
%$BM generalizes 
$0/1${\sc-knapsack} admits a Fully PTAS, it is natural to explore the existence of a Fully PTAS for BM. There are known Fully PTASs  for BM on restricted families of matroids. This includes {\sc knapsack with a cardinality constraint}~\cite{caprara2000approximation}, 
{\sc multiple-choice knapsack}~\cite{La79}, and {\sc BM with laminar matroid constraint} \cite{DKS23c}. However, the question of whether BM admits a Fully PTAS on general matroids remained open. 

In this paper, we resolve this question negatively for BM and other fundamental matroid optimization problems with a linear constraint. 

\subsection{Our Results}

For a matroid $\cm = (E,\cI)$ we define $\IS(\cm) = \cI$ and $\bases(\cm) = \{S \in \cI~|~|S| = \textsf{rank}(\cm)\}$, where $\textsf{rank}(\cm) = \max_{T \in \cI} |T|$ is the {\em rank} of $\cm$, i.e., the maximum cardinality of an independent set. We study a family of {\em matroid optimization problems with a linear constraint} (MOL). Problems in this family are characterized by three parameters: 
\begin{enumerate}
\item[(i)]
%$(i)$ 
The optimization objective $\opt$ $-$ either the operator $``\max"$ or $``\min"$. 
\item[(ii)]
%$(ii)$ 
A {\em matroid feasibility constraint} $\cF$ $-$ either the independent sets of a matroid, or the set of {\em bases} of a matroid. The feasibility constraint is $\cF \in \{\IS,\bases\}$.
\item[(iii)]
%$(iii)$ 
A relation $\triangleleft$ $-$  realized by one of the relations $``\geq"$ or $``\leq"$. 
\end{enumerate}

Let $\cP = \{\max,\min\} \times \{\bases,\IS\} \times \{\leq ,\geq\}$ be the set of {\em parameters} for MOL problems. 
Based on the set of parameters $\cP$, we define for every triplet a problem
in the MOL family. For $P \in \cP$ where $P = \left(\opt,\cF,\triangleleft\right)$, define the $P$-{\sc matroid optimization with a linear constraint} ($P$-MOL) problem as follows. An instance is a tuple $I = (E,\cI,v,w,L)$ such that $\cm = (E,\cI)$ is a matroid, $v:E \rightarrow \mathbb{R}_{\geq 0}$ is the objective function, $w:E \rightarrow \mathbb{R}_{\geq 0}$ is a weight function, and $L \in \mathbb{R}_{\geq 0}$ is a {\em bound} for the linear constraint. 
A {\em solution} of $I$ is $S \subseteq E$ which satisfies the matroid feasibility constraint $S \in \cF(\cm)$ and the linear constraint $w(S) \triangleleft L$. The goal is to optimize (i.e., maximize or minimize) the value $v(S)$. Thus, we can formulate a P-MOL optimization problem as
\begin{equation}
\label{eq:0}
\opt~ v(S) \text{ s.t. } S \in \cF(\cm),~  w(S) \triangleleft L.
\end{equation} 



Observe that $(\max,\IS,\leq)$-MOL is the BM problem.
That is, given a BM instance (equivalently, $(\max,\IS,\leq)$-MOL instance) $I = (E,\cI,v,w,L)$, the goal is to find an independent set $S \in \cI$ of maximum total value $v(S)$ such that $w(S) \leq L$. Other notable examples for MOL problems are {\sc constrained minimum basis of a matroid} (CMB) \cite{HL04}, which can be cast as $(\min,\bases,\leq)$-MOL, and {\sc knapsack cover with matroid constraint} (KCM) \cite{CCNR13} formalized by $(\min,\IS,\geq)$-MOL.\footnote{CMB, KCM, and other MOL problems may not have a solution; however, we can decide in polynomial time if a solution exists, and our definition of approximation algorithms captures instances with no solution.}



We note that \eqref{eq:0} does not refer to the {\em representation} of the instance $I$. We consider two possible representations.
For any $P \in \cP$, in an instance $(E,\cI,v,w,L)$ of {\em oracle} $P$-MOL, the arguments
$E,v,w,L$ are given as the input, and the independent sets $\cI$ are accessed via a {\em membership oracle}, which determines whether a given set $S \subseteq E$ belongs to  $\cI$ in a single query. Thus, the independent sets are not considered in the encoding size of the instance. The term {\em running time} for problems involving oracles refers to the sum of the number of queries to the oracle and the number of basic operations. 
Previous works on MOL problems often consider membership oracles \cite{CVZ11,BBGS11,CCNR13,DKS23,DKS23b}. As hardness with oracles does not necessarily imply hardness in non-oracle models (see, e.g., \cite{canetti2004random,chang1994random}), in \Cref{sec:SAT} we show lower bounds for variants of MOL problems in which the independent sets are encoded as part of the input.   




Clearly, the problem $(\min,\IS,\leq)$-MOL is trivial since the empty set achieves the optimal objective value. 
However, for any other $P \in \cP$, solving the $P$-MOL problem is challenging.
The {\em non-trivial} MOL problems are all the MOL problems excluding $(\min,\IS,\leq)$-MOL. That is, $P$-MOL is non-trivial if $P\in \cQ$ where  $\cQ = \cP \setminus \{(\min,\IS,\leq)\}$.
%We consider the {\em non-trivial} MOL problems, whose parameters exclude the above triplet, i.e., $\cQ = \cP \setminus \{(\min,\IS,\leq)\}$.
 Observe that non-trivial MOL problems are NP-hard (e.g., $0/1${\sc-knapsack} is a special case of $(\max,\IS,\leq)$-MOL);
however, %several 
all previously studied MOL problems admit approximation schemes.




For certain special cases of MOL problems, e.g., BM with simple matroid constraints,
the existence of a Fully PTAS is known for decades \cite{Si79,caprara2000approximation}.
However, for MOL problems with arbitrary matroid constraints, the best known results are Efficient PTAS.  While matroids form an important generalization of well known basic
constraints,
% (cardinality constraint, linear independence in vector space, cycle-free subsets of edges in a graph, etc.), 
the complexity of the corresponding MOL problems remained open. Specifically, prior to this work, the existence of a MOL problem which does not admit a Fully PTAS was open.





Our main result is that none of the (non-trivial) 
oracle matroid optimization problem with a linear constraint admits 
a Fully PTAS, even if randomization is allowed. This unconditioned hardness result is established by deriving a lower bound on the minimum number of queries to the membership oracle. 
 \begin{restatable}{theorem}{thmMain}
	\label{thm:main}
	For every $P \in \cQ$ there is no randomized  \textnormal{Fully PTAS} for oracle $P$-\textnormal{MOL}. 
\end{restatable} 

\begin{table}[tbh]
	\centering
	\begin{tabular}{m{2.7in} m{1.4in} m{1.9in}}
		\toprule
		Problem & Previous Results & This Paper \\
		\midrule
		Budgeted Matroid Independent Set & Efficient PTAS \cite{DKS23} & {\bf No Fully PTAS} \\
		Budgeted Matroid Intersection & Efficient PTAS  \cite{DKS23b} & {\bf No Fully PTAS}  \\
		Constrained Minimum Basis of Matroid & Efficient PTAS  \cite{HL04} & {\bf No Fully PTAS} \\
		Knapsack Cover with a Matroid & PTAS \cite{CCNR13} & {\bf No Fully PTAS}  \\
		\bottomrule
	\end{tabular}
	\caption{Implications of our results for previously studied MOL problems. All of our bounds hold for randomized algorithms. }
	\label{tab:results}
\end{table}

\Cref{thm:main} conclusively distinguishes MOL problems with arbitrary matroids, such as BM, from special cases with simpler matroid constraints. Furthermore, it shows that existing Efficient PTAS \cite{HL04,DKS23,DKS23b} for MOL problems on general matroids are the best possible. Notable implications of our results are given in \Cref{tab:results}, and consequences of our lower bounds for a set of previously studied problems~\cite{ CGM92,HL04,BBGS11,CVZ11,GZ10,CCNR13,DKS23,DKS23b,doron2023budgeted}
are given in \Cref{sec:results}. By resolving the complexity status of MOL problems on general matroids, our results promote future research to design (or rule out) Fully PTAS for MOL problems on restricted matroid classes (see \Cref{sec:discussion}). 

To prove \Cref{thm:main}, we turn our attention to the following problem. %- a generalization of {\sc subset sum} with a matroid constraint. 
\begin{definition}
	\label{def:EMB}
	An instance of {\sc Exact Matroid Basis (EMB)}  is $I =  (E,\cI,c,T)$, where $(E,\cI)$ is a matroid, $c:E \rightarrow \mathbb{N}$ is a weight function, and $T \in \mathbb{N}$ is a target value. A {\em solution} is a basis $S$ of $(E,\cI)$ such that $c(S)=T$. The goal is to decide if there is a solution. 
\end{definition}
Similar to MOL problems, EMB does not specify the input. 
In an instance $(E,\cI,c,T)$ of {\sc oracle-EMB},
$E,c,T$ are explicitly given, and the independent sets $\cI$ are accessed via a membership oracle. An instance $I$ of a decision problem $\cD$
is a ``yes"-instance if the correct answer for $I$ is ``yes"; otherwise, $I$ is a ``no"-instance.
We say that $\cA$ is a {\em randomized algorithm} for a decision problem $\cD$ if,
% A {\em randomized algorithm} $\cA$ for a decision problem $\cD$ is an algorithm that,
 given a ``yes"-instance $I$ of $\cD$, $\cA$ returns ``yes"  with probability at least $\frac{1}{2}$; for a ``no"-instance, $\cA$ returns ``no" with probability $1$. The next result rules out a pseudo-polynomial time algorithm for oracle-EMB, thus distinguishing the problem from the special cases of $k$-{\sc subset sum} and EMB on linear matroids, which admit a pseudo-polynomial  time algorithm \cite{CGM92}. 
% We prove that the classic (randomized) pseudo-polynomial algorithm for oracle-EMB on {\em representable} (or, linear) matroids \cite{CGM92}, cannot be extended to general matroids.  

\begin{restatable}{theorem}{thmemb}
	\label{thm:emb}
	There is no randomized  algorithm  for {\sc oracle-EMB} that for any  \textnormal{oracle-EMB} instance $I =(E,\cI,c,T)$, where $n = |E|+1$ and  $m = c(E)+1$,  runs in time $ \left(n \cdot(T+2) \cdot m \right)^{O(1)}$. 
\end{restatable}

 

\subsection{Technical Contribution}
We derive our results by introducing $\Pi$-{\em matroids}.
This new family of matroids carefully exploits a simple weight function to define a matroid that successfully {\em hides} a specific property $\Pi$ within its independent sets (see \Cref{sec:technical}). Using $\Pi$-matroids, we define %the $\Pi$-{\em matroid problem} and prove its hardness; 
	oracle-EMB instances whose solutions must satisfy the property $\Pi$. This shows the unconditional hardness of oracle-EMB, as $\Pi$ can be discovered only via an exponential number of queries to the membership oracle. Our hardness results for MOL problems (as stated in \Cref{thm:main}) are derived via reduction from oracle-EMB.
	
	
		Despite the abundance of lower bounds for matroid problems~% , for example, the powerful bounds of Jensen and Korte~
		\cite{lovasz1978matroid,JK82,karp1985complexity,soto2014simple}, as well as for knapsack  problems~\cite{chekuri2005polynomial,kulik2010there,cygan2019problems,bringmann2021fine}, we are not aware of lower bounds that leverage the interaction between the matroid constraint and the additional linear constraint required for deriving our new lower bound for EMB, and consequently for MOL problems. Indeed, if the matroid constraint is removed from \eqref{eq:0} (equivalently, $\cF(\cm)=2^E$), MOL problems become variants of classic 
	$0/1${\sc-knapsack}, which admits a Fully PTAS. Alternatively, if the linear constraint imposed by $w,L$ is removed, then we have the polynomially solvable maximum/minimum weight matroid independent set problem. This distinguishes our construction from existing lower bounds for matroid problems. %We believe that
	 $\Pi$-matroids may be useful for deriving lower bounds for other problems %, especially for problems whose hardness comes from the interaction between a matroid constraint and a non-uniform weight function 
	(see \Cref{sec:discussion}).  
	

	
	
	
		 Our unconditional lower bounds apply in the {\em oracle} model, where the independent sets of the given matroid can be accessed only via a membership oracle. One may question the validity of the bounds for variants of the problems where the matroid is encoded as part of the input. Indeed, in some scenarios, the use of oracles makes problems harder \cite{canetti2004random,chang1994random}. Thus, we complement our results by showing that the same lower bounds hold under the standard complexity assumption $\textnormal{P}\neq\textnormal{NP}$, even if the matroid is encoded as part of the instance and membership can be decided in polynomial time. We accomplish this by designing the family of {\em \textnormal{SAT}-matroids} $-$ a counterpart of the $\Pi$-matroid family whose members can be efficiently encoded. This construction can be used to obtain hardness results for other matroid problems in non-oracle models, based on existing analogous lower bounds in the oracle model (e.g., \cite{JK82}). We elaborate on that in \Cref{sec:SAT}.
		 %This construction may be useful in deriving hardness results for matroid problems with no membership oracle. We elaborate on that in \Cref{sec:SAT}.
	
	

